\section{Method}
\label{sec:method}

\subsection{Overview}
\label{sec:overview}

UrbanDreamer learns cooperative routing policies under unknown future demand
by reusing environment experience through a learned traffic world model.
As illustrated in Fig.~\ref{fig:architecture}, the framework couples
network-wide traffic prediction with vehicle-level route selection:
the world model captures how routing decisions change road-segment traffic,
while the shared actor selects remaining routes using vehicle-specific states
and available traffic information. Following
Section~\ref{sec:wm-formulation}, vehicle decisions and network-wide traffic
predictions follow the same fixed physical-time grid.

After representation learning and world-model pretraining on archived
traffic trajectories, online training repeats four stages. First, the
current actor interacts with SUMO, and the observed traffic, vehicle states,
realized requests, executed route choices, turning commitments, and global
rewards are stored in a replay buffer. Second, the world model is fitted
to recorded transitions and their corresponding controls. Third, imagined
rollouts start from replay states and the associated active vehicle records.
The actor generates new route choices at fixed decision times, while the
learned traffic dynamics and vehicle execution produce subsequent imagined
states. These rollouts advance vehicles already active at initialization.
Finally, centralized value learning uses return targets from these
rollouts, and global reward-to-go advantages guide updates to the shared
actor. Repeating this cycle supplies additional policy-training
trajectories without requiring a new external rollout for each imagined
trajectory.

The traffic model follows a coarse-to-fine design: a low-frequency graph
prior provides a coarse estimate of the latent state increment, and
conditional flow matching uses this estimate as a background condition
when generating the next state from the current state with graph-structured
noise. To support recursive use, model fitting includes multi-step
supervision, while a latent-state disagreement penalty regularizes the
global imagined reward used for policy optimization. We describe the causal
routing framework in Section~\ref{sec:wm-routing}, the graph-structured
traffic model in Section~\ref{sec:traffic-world-model}, and the stability
mechanisms and actor--critic updates in
Section~\ref{sec:stable-imagination}.

\begin{figure*}[t]
  \centering
  \includegraphics[width=0.8\textwidth]{figure/UrbanDreamer}
  \caption{Overview of UrbanDreamer. Environment experience trains the
  graph-structured traffic model; imagined rollouts support global value
  learning and cooperative routing policy optimization.}
  \label{fig:architecture}
\end{figure*}

\subsection{World-Model-Based Cooperative Routing}
\label{sec:wm-routing}

\subsubsection{Experience Collection and Imagination Initialization}

UrbanDreamer converts collected traffic experience into new routing
trajectories through model-based imagination. Replay $\mathcal B$ stores
macro-grid traffic fields, explicit vehicle records, and their matching
requests, controls, and rewards. Contiguous sequences train the traffic
model. To initialize imagination, we sample uniformly over eligible
nonterminal macro boundaries containing active vehicles and restore their
positions, destinations, remaining routes, and movement progress. The
observed traffic field is encoded with the current encoder. The original
physical time, cumulative travel cost $C_{\mathrm{time}}$, and trip-start
count $B$ are also restored. From this restored state, the actor generates
new decisions and the model supplies subsequent traffic predictions,
extending the training experience beyond the recorded prefix.

\subsubsection{Periodic Causal Imagination}
\label{sec:causal-imagination}

At decision time $s_k$, the shared actor scores candidate remaining routes
for each $i\in\mathcal D_k$:
\begin{equation}
 \pi_\theta(a\mid o_{i,k})=
 \frac{\exp f_\theta^{\mathrm{route}}(o_{i,k},a)}
 {\sum_{b\in\mathcal A_i(S_k)}
       \exp f_\theta^{\mathrm{route}}(o_{i,k},b)}.
 \label{eq:shared-next-hop-policy}
\end{equation}
Vehicles follow the selected routes between decision times. Each imagined
environment starts with the vehicles already departed and not yet arrived
at its sampled boundary. The restored cumulative trip-start count remains
fixed during the rollout, and future requests are not inserted. Real environment collection
continues to reveal requests online; the imagined vehicle layer therefore
approximates subsequent evolution of the observed active cohort.

Macro-endpoint updates couple periodic route selection with traffic
prediction. At $t_n$, eligible vehicles select routes from a shared
pre-action traffic snapshot. Within $[t_n,t_{n+1})$, the vehicle layer
advances using traffic-speed readouts from $\widehat Z_n$, held fixed
within that interval, and local execution rules. It processes arrivals
and records the turning commitments encountered while following the
selected routes. At $t_{n+1}$, this completed control log forms $\Psi_n$.
The world model uses the current latent state and these controls to
predict the interval's endpoint; its readouts then drive the next interval.
Boundary route choices and turning commitments belong to the next
half-open interval. This split update handles local execution within
each interval and refreshes the remaining network-level traffic feedback
at its boundary.

Explicit records maintain vehicle identities, trip starts, destination
arrivals, and the counts used for reward accounting. The decoder supplies
road-segment count and speed readouts. At each fixed state and condition,
$K$ independent latent samples provide the disagreement statistic.
We propagate one designated sample,
$\widehat Z_{n+1}=\widehat Z_{n+1}^{(1)}$.
Its selection is independent of the sample values, so it retains the
model's conditional transition distribution and cross-segment dependencies.
All $K$ samples contribute to disagreement estimation. An imagined rollout
lasts at most eight 30-second intervals and ends earlier at the physical
cutoff or when all vehicles in its initial active cohort have arrived.

\subsection{Graph-Structured Traffic World Model}
\label{sec:traffic-world-model}

We organize traffic transition learning into spatially coherent change
and conditional refinement. A graph prior predicts a low-frequency
background increment and supplies it as an explicit transition condition.
Conditional flow matching transforms the current state with
graph-structured noise into complete next-state samples, capturing
nonlinear variation and accommodating errors in the coarse estimate.
This coarse-to-fine design uses road structure in both state representation
and transition generation. Compared with propagating a single point
prediction, conditional sampling supplies alternative network states for
imagined policy learning, exposing the actor to different downstream
congestion consequences under the available state and controls.

\subsubsection{Graph Representation and Transition Conditions}

A residual graph encoder with layer normalization aggregates traffic
along the directed road graph into
$Z_n=f_\phi(X_n,G)\in\mathbb R^{M\times d}$.
Its segment-aligned features combine local conditions with neighboring
traffic. The decoder gives $\widehat Y_n=d_\phi(Z_n)$, including vehicle
counts and speeds, together with occupancy and normalized halting counts;
static attributes remain fixed. Counts are reconstructed in vehicle units.
For speed supervision, the target is the mean of the observed speeds at
the current and following boundaries, normalized by the segment speed
limit. Thus the decoded speed at $t_n$ represents the following interval.

Let $\langle\cdot\rangle$ denote a mean over valid segment observations in
a reconstruction batch, with subscripts $\mathrm{occ}$ and
$\mathrm{slow}$ restricting it to occupied and sufficiently populated
slow segments. A mean over an empty subset is zero. Define
$e_u^c=[\log(1+\widehat c_u)-\log(1+c_u)]^2$ and let $e_u^v$ be the
smooth-$L_1$ loss, with threshold $0.1$, between the normalized predicted
speed and its forward-interval target. Reconstruction uses
\begin{equation}
 \begin{aligned}
 \mathcal L_c
   &=(1-\zeta)\langle e^c\rangle
      +\zeta\langle e^c\rangle_{\mathrm{occ}},\\
 \mathcal L_v
   &=\langle e^v\rangle
      +w_{\mathrm{occ}}\langle e^v\rangle_{\mathrm{occ}}
      +w_{\mathrm{slow}}\langle e^v\rangle_{\mathrm{slow}},\\
 \mathcal L_{\mathrm{mass}}
   &=\mathbb E_{\mathcal B}\!\left[
      \operatorname{asinh}^{2}\!\left(
       \frac{\sum_u(\widehat c_u-c_u)}{1+\sum_u c_u}
      \right)\right],\\
 \mathcal L_{\mathrm{rec}}
   &=\mathcal L_c+0.1\mathcal L_{\mathrm{mass}}+\mathcal L_v
       +0.05(\mathcal L_{\mathrm{occ}}+\mathcal L_{\mathrm{halt}}).
 \end{aligned}
 \label{eq:traffic-reconstruction}
\end{equation}
Here $\zeta=0.2$, set to zero for a batch without occupied segments;
$\mathcal L_{\mathrm{occ}}$ and $\mathcal L_{\mathrm{halt}}$ are mean
squared errors in occupancy and normalized halting counts. The speed
weights and slow-segment thresholds are configured using the training
protocol. Speed targets require an observed following boundary.
After reconstruction pretraining, the encoder is frozen during dynamics
learning. Decoder calibration uses a separate supervised objective and
optimizer, including during online learning. Neither encoder nor decoder
is updated through the flow-matching loss.

The transition condition is $C_n=\Psi_n$. It encodes the observed or
imagined turning commitments within the interval, including source
segment, next segment, destination, and relative occurrence time.
The transition is
\begin{equation}
 \widehat Z_{n+1}=F_\phi(Z_n,C_n;\eta_n).
 \label{eq:causal-macro-history}
\end{equation}
The velocity network processes the current physical transition as one
traffic block, using segment tokens, a flow-time token, and register
tokens. Its segment tokens receive the encoded controls and a projection
of the current latent state. Recursive predictions become the current
state for subsequent transitions. Clean future states supply fitting
targets rather than transition conditions.

\subsubsection{Low-Frequency Graph Prior}

The directed adjacency $A$ encodes legal vehicle movements and directed
message passing. To extract spatially coherent changes, we construct
the symmetric weights $W=(A+A^\top)/2$. With
$D_W=\operatorname{diag}(W\mathbf1)$, use the symmetric normalized
Laplacian and its soft low-pass operator:
\begin{equation}
 \begin{aligned}
 L_G&=D_W^{\dagger/2}(D_W-W)D_W^{\dagger/2},\\
 P_{\mathrm{soft}}&=(I+\tau_{\mathrm{LP}}L_G)^{-1},
 \qquad \tau_{\mathrm{LP}}\geq0.
 \end{aligned}
 \label{eq:prior-spectral-operator}
\end{equation}
The diagonal entries of $D_W^{\dagger/2}$ are inverse square roots of
positive degrees and zero for isolated segments. Low graph frequencies
describe spatially coherent variation~\cite{shuman2013emerging}.
For $L_G=U\Lambda U^\top$, the filter has spectral response
\begin{equation}
 P_{\mathrm{soft}}
   =U\operatorname{diag}\!\left(
      \frac{1}{1+\tau_{\mathrm{LP}}\lambda_j}
     \right)U^\top.
 \label{eq:soft-low-frequency-filter}
\end{equation}
It attenuates higher graph frequencies while retaining all modes.
The mainline uses $\tau_{\mathrm{LP}}=0.1$.

With $\Delta Z_n=Z_{n+1}^{\mathrm{tar}}-Z_n$, the state-only prior is
\begin{equation}
 \begin{aligned}
 R_n&=p_\phi(Z_n;A),\\
 b_n=\Delta\widehat Z_n^{\mathrm{low}}
    &=P_{\mathrm{soft}}R_n,\\
 \mathcal L_{\mathrm{prior}}
   &=\mathbb E_{\mathcal B}\!\left[
       \frac{1}{Md}
       \left\|b_n-\operatorname{sg}
          (P_{\mathrm{soft}}\Delta Z_n)\right\|_F^2
      \right].
 \end{aligned}
 \label{eq:graph-prior-loss}
\end{equation}
Here $\operatorname{sg}$ denotes a stop-gradient operation.
The prior combines the current state with separate upstream and downstream
messages. Its input retains all graph frequencies, and filtering is
applied only to the predicted increment. The background increment $b_n$
is supplied to the velocity network with stopped gradients; the auxiliary
prior loss trains the background predictor.

Graph-structured noise defines the transition source around the current
state:
\begin{equation}
 \begin{aligned}
 \bar x_{0,n}&=Z_n,\\
 \tilde x_{0,n}&=Z_n+\sigma U_r\eta_n,
 \qquad (\eta_n)_{ab}\overset{\mathrm{iid}}{\sim}\mathcal N(0,1),
 \end{aligned}
 \label{eq:noisy-graph-prior}
\end{equation}
where $\eta_n\in\mathbb R^{r\times d}$ and $U_r$ is an orthonormal
basis of $r$ lowest-frequency graph modes. The noise basis is distinct
from the full-rank soft filter used by the background predictor.
The mainline uses $\sigma=0.1$ and rank $r=16$ where the graph size permits.
The current state remains the source mean and reference state, while
$b_n$ provides the coarse predictive condition. The learned transport
operates in the complete latent space, allowing both spatially coherent
and localized variation.

\subsubsection{Conditional Flow-Matching Transition}

The conditional velocity field transports the noisy current state to the
next latent state. For a recorded pair, sample
$\tau\sim\mathcal U(0,1)$ and define
\begin{equation}
 \begin{aligned}
 x_{n,\tau}&=(1-\tau)\tilde x_{0,n}
                  +\tau Z_{n+1}^{\mathrm{tar}},\\
 \nu_n^\star&=Z_{n+1}^{\mathrm{tar}}-\tilde x_{0,n},\\
 Z_{n+1}^{\mathrm{tar}}
   &=\operatorname{sg}[f_\phi(X_{n+1}^{\mathrm{tar}},G)].
 \end{aligned}
 \label{eq:conditional-fm-path}
\end{equation}
The target field $X_{n+1}^{\mathrm{tar}}$ uses the vehicle-membership
convention specified in the multi-step supervision below.
The velocity network predicts
$\widehat\nu_{n,\tau}
 =\nu_\phi(x_{n,\tau},\tau,C_n,Z_n,\operatorname{sg}(b_n))$,
trained by conditional flow matching~\cite{lipman2023flow}:
\begin{equation}
 \mathcal L_{\mathrm{FM}}
   =\mathbb E_{\mathcal B,\tau,\eta}\!\left[
       \frac1{Md}
       \left\|\widehat\nu_{n,\tau}
                  -\operatorname{sg}(\nu_n^\star)\right\|_F^2
     \right].
 \label{eq:conditional-fm-loss}
\end{equation}
The loss weights latent coordinates uniformly. In addition to the
transformer features, a conditional residual head receives the current
state, the detached background increment, the displacement
$x_{n,\tau}-Z_n$, and flow time. The encoder supplies detached reference
targets. During recursive training, gradients remain live through the
generated current state and its reference projection, while the background
condition remains detached.

Inference integrates from $\hat x_n(0)=\tilde x_{0,n}$:
\begin{equation}
 \begin{aligned}
 \frac{\mathrm d\hat x_n(\tau)}{\mathrm d\tau}
   &=\nu_\phi\bigl(\hat x_n(\tau),\tau,C_n,Z_n,
                         \operatorname{sg}(b_n)\bigr),\\
 \widehat Z_{n+1}&=\hat x_n(1).
 \end{aligned}
 \label{eq:conditional-fm-ode}
\end{equation}
Euler integration operates in the full latent space. The parameter $\tau$
tracks transport from the source to the endpoint, and each completed solve
predicts one physical interval $\Delta T$. The mainline uses 100 Euler
steps per transition. The basic dynamics objective combines prior
alignment and flow matching:
\begin{equation}
 \mathcal L_{\mathrm{base}}
   =\mathcal L_{\mathrm{FM}}
      +\alpha_{\mathrm{prior}}\mathcal L_{\mathrm{prior}},
 \qquad \alpha_{\mathrm{prior}}=0.1.
 \label{eq:base-world-model-loss}
\end{equation}
Reconstruction and decoder calibration are optimized separately from
this latent dynamics objective. An auxiliary exposure readout is also
fitted by squared error using detached observed endpoint latents and
recorded active-vehicle exposure; imagined rewards use the explicit
vehicle accounting described below.

\subsection{Stable Imagined Policy Learning}
\label{sec:stable-imagination}

We couple model-side trajectory supervision with policy-side reward
regularization. Multi-step supervision trains the traffic model on its own
recursive predictions. Latent disagreement assigns a penalty to
high-dispersion predictions during policy learning.

\subsubsection{Multi-Step Model Supervision}

To align model fitting with recursive imagination, initialize a contiguous
replay sequence at
$\widehat Z_{n\mid n}=\operatorname{sg}[f_\phi(X_n,G)]$.
Let $\mathcal C_n^0$ contain the vehicles already active at the sequence
start. At later recorded boundaries, vehicle-count targets include only
members of $\mathcal C_n^0$ that remain on the network; speed, occupancy,
and halting labels retain their observed background values. Recorded
turning controls are restricted to the same cohort. Denote the resulting
targets and controls by $X_{n+h}^{\mathrm{tar}}$ and
$C_{n+h-1}^{\mathcal B,\mathcal C_n^0}$, respectively. Predict
\begin{equation}
 \widehat Z_{n+h\mid n}
   =F_\phi\bigl(\widehat Z_{n+h-1\mid n},
          C_{n+h-1}^{\mathcal B,\mathcal C_n^0};\eta_{n+h-1}\bigr).
 \label{eq:multistep-model-rollout}
\end{equation}
Each sampled prediction becomes the next input. At every depth, construct
the noisy source and interpolated flow path from this generated input
and the corresponding detached recorded target, using
Eqs.~\eqref{eq:noisy-graph-prior}--\eqref{eq:conditional-fm-loss}.
Let $\ell_{n,h}^{\mathrm{FM}}$ and
$\ell_{n,h}^{\mathrm{prior}}$ be the resulting velocity-regression and
prior-alignment losses at depth $h$. The supervised discrepancy is
\begin{equation}
 \ell_{n,h}^{\mathrm{roll}}
   =\ell_{n,h}^{\mathrm{FM}}
       +\alpha_{\mathrm{prior}}\ell_{n,h}^{\mathrm{prior}}.
 \label{eq:explicit-rollout-discrepancy}
\end{equation}
If $b_{n,h}$ marks valid recorded endpoints and
$H_n^{\mathrm{valid}}=\sum_h b_{n,h}>0$, the recursive objective is
\begin{equation}
 \mathcal L_{\mathrm{wm}}
   =\mathbb E_{\mathcal B,\tau,\eta}\!\left[
       \frac{\sum_{h=1}^{H_{\mathrm{wm}}}
                       b_{n,h}\ell_{n,h}^{\mathrm{roll}}}
            {H_n^{\mathrm{valid}}}
      \right].
 \label{eq:full-world-model-loss}
\end{equation}
The default supervision depth is $H_{\mathrm{wm}}=3$, corresponding to
90 seconds; shorter sequences use their available recorded endpoints.
The policy horizon $H_\pi=8$ and solver resolution are configured
separately. Gradients propagate through the finite predicted unroll,
including earlier flow solves, while recorded targets and target
velocities are detached. Thus multi-step supervision applies the
flow-matching objective at recursively generated prefixes. The
disagreement statistic is used for imagined-reward regularization rather
than for weighting this model-fitting objective.

\subsubsection{Latent-State Disagreement Regularization}
\label{sec:global-return}

Holding $(Z,C)$ fixed, generate
$\widehat Z^{(j)}=F_\phi(Z,C;\eta^{(j)})$, $j=1,\ldots,K$,
with $K\ge2$. Using the raw latent coordinates, define
\begin{equation}
 \begin{aligned}
 \overline Z&=\frac1K\sum_j\widehat Z^{(j)},\\
 d_u(Z,C)&=\frac{\sum_j
       \|\widehat z_u^{(j)}-\overline z_u\|_2^2}{(K-1)d},\\
 u(Z,C)&=\frac1M\sum_u d_u(Z,C).
 \end{aligned}
 \label{eq:latent-disagreement}
\end{equation}
This statistic measures conditional prediction dispersion directly in
latent space, with identity scaling of latent coordinates. Its network
average supplies the global imagined-reward penalty. It measures
dispersion under the learned conditional model, rather than an error bound
or a calibrated estimate of epistemic uncertainty.

Policy learning uses the finite-horizon cumulative mean
$m(t)=C_{\mathrm{time}}(t)/(B(t)\vee1)$ defined in the preliminaries.
All started trips contribute their elapsed travel time up to arrival or
the cutoff $T$. For completeness, a discount-compatible form follows by
setting $\Gamma(a,b)=\gamma_0^{(b-a)/\Delta T}$ and writing
\begin{equation}
 r_\gamma^g[a,b]=\kappa_g[m(a)-\Gamma(a,b)m(b)].
 \label{eq:method-global-reward}
\end{equation}
Across $[0,T]$, its discounted sum telescopes to
$-\kappa_g\Gamma(0,T)m(T)$, retaining the fixed-horizon mean objective.
The implemented learner uses $\gamma_0=1$ and $\kappa_g=1/60$, recovering
the undiscounted mean-increment reward. During imagination, $B(t)$ remains
equal to the restored trip-start count, and $C_{\mathrm{time}}$ grows by
the explicitly accumulated travel time of the active seed vehicles.

Let $q_j$ denote successive imagined macro boundaries, including a
possibly shorter terminal endpoint, and write
$u_j=u(\widehat Z_j,C_j)$. Charge the penalty once for each complete
predicted physical interval:
\begin{equation}
 \widetilde r_j^g
   =r^g[q_j,q_{j+1}]
     -\beta\,\mathbf1\{q_{j+1}-q_j=\Delta T\}\,u_j.
 \label{eq:regularized-event-reward}
\end{equation}
The mainline uses $K=2$ and $\beta=0.1$. A partial terminal interval
contributes its observed physical cost without an additional model
transition or disagreement charge. The actor--critic learner optimizes
these regularized returns; traffic evaluation uses the finite-horizon
mean travel-time objective.

\subsubsection{Global Actor--Critic Optimization}

A centralized state-value critic $V_\omega(\widehat Z_j,c_j)$ estimates
the shared remaining return at $q_j$. It receives mean-pooled segment
latents and the causal context
$c_j=[q_j/T,\,N_{\mathrm{active}}(q_j)/(B(q_j)\vee1),\,
\log(1+C_{\mathrm{time}}(q_j)/60),\,\log(1+B(q_j))]$.
Let $v_j^-=\operatorname{sg}[V_\omega(\widehat Z_j,c_j)]$,
setting it to zero at the physical cutoff or exhaustion of the imagined
seed cohort. Let $d_j$ mark either terminal condition at $q_{j+1}$.
For a rollout ending at $q_L$, use
\begin{equation}
 \begin{aligned}
 G_L^\lambda&=v_L^-,\\
 G_j^\lambda&=\widetilde r_j^g+
       (1-d_j)[(1-\lambda)v_{j+1}^-+\lambda G_{j+1}^\lambda],\\
 \mathcal L_V&=\mathbb E_j\!\left[
       \bigl(V_\omega(\widehat Z_j,c_j)
                -\operatorname{sg}(G_j^\lambda)\bigr)^2\right],
 \end{aligned}
 \label{eq:global-lambda-return}
\end{equation}
with $\lambda=0.95$. At a nonterminal imagination-horizon cutoff, the
critic supplies the bootstrap value. It continues to value shared
outcomes after individual vehicle arrivals while other cohort vehicles
remain active.

A shared candidate-action critic $Q_\psi$ supplies a vehicle-specific
counterfactual baseline under the same global reward. Its causal input
$\chi_{i,k}$ includes the current traffic latent, route-candidate
features, destination, reward context, and other simultaneous choices.
For $s_k=q_{j(k)}$, the decisions share a Monte Carlo reward-to-go target
with value bootstrapping at a nonterminal rollout cutoff:
\begin{equation}
 \begin{aligned}
 R_k&=\sum_{j=j(k)}^{L-1}\widetilde r_j^g+v_L^-,\\
 b_{i,k}&=\operatorname{sg}\!\left[
    \sum_{a\in\mathcal A_i(S_k)}
       \pi_\theta(a\mid o_{i,k})Q_\psi(\chi_{i,k},a)\right],\\
 A_{i,k}&=\operatorname{sg}(R_k-b_{i,k}),\\
 \mathcal L_Q&=\mathbb E_{i,k}\!\left[
      \bigl(Q_\psi(\chi_{i,k},a_{i,k})
                   -\operatorname{sg}(R_k)\bigr)^2\right].
 \end{aligned}
 \label{eq:counterfactual-global-advantage}
\end{equation}
The baseline is saved before fitting either critic and reused for the
actor update. Only executed candidate actions are fitted to $R_k$.
Simultaneous decisions share their global return target, while their
counterfactual baselines can differ.

Across the decision samples in an imagination training pass, normalize
advantages by their population standard deviation, floored at one,
without subtracting their mean:
\begin{equation}
 \begin{aligned}
 s_A&=\max\{1,\operatorname{Std}(\{A_{i,k}\})\},\\
 \widehat A_{i,k}&=\operatorname{clip}(A_{i,k}/s_A,-2.5,2.5).
 \end{aligned}
 \label{eq:global-advantage-scaling}
\end{equation}
The shared actor uses the entropy-regularized policy-gradient objective
\begin{equation}
 \begin{aligned}
 \mathcal L_\pi=-\mathbb E_{i,k}\bigl[
    &\log\pi_\theta(a_{i,k}\mid o_{i,k})
                 \operatorname{sg}(\widehat A_{i,k})\\
    &+\alpha_{\mathrm{ent}}\mathcal H(
                 \pi_\theta(\cdot\mid o_{i,k}))\bigr].
 \end{aligned}
 \label{eq:next-hop-actor-loss}
\end{equation}
The expectation averages over vehicle decisions, and $\mathcal H(\pi)$
denotes policy entropy. The entropy coefficient follows the configured
episode-based annealing schedule. Each periodic route decision
contributes a policy update weighted by its global reward-to-go advantage.
Actor--critic optimization treats imagined trajectories and saved
baselines as fixed training data, with model parameters held constant.
Alternating this optimization with environment collection and model
fitting converts traffic experience into additional cooperative-policy
updates.
